% !TeX root = ../抽象代数笔记.tex
% 由主文档通过 \input 载入；请编译主文档。
\setcounter{chapter}{17} % 保留原书章号。
\chapter{伽罗瓦对应与伽罗瓦群计算}\label{ch:galois}

\begin{description}
  \item[学习目标] 在子群与中间域之间双向转换，并从多项式根的置换计算具体扩张的对称性。
  \item[先修知识] 第\ref{ch:quotient-groups}、\ref{ch:actions}、\ref{ch:splitting}章；有限域提供一类完整例子。
\end{description}

\section{域自同构、固定域与伽罗瓦群}
对扩张 $L/F$ 定义固定 $F$ 的自同构群 $\operatorname{Aut}_F(L)$；在有限伽罗瓦情形记为 $\operatorname{Gal}(L/F)$。对自同构子群 $H$ 定义固定域 $L^H=\{a\in L:\sigma(a)=a\text{ 对所有 }\sigma\in H\}$。从自同构对根的作用说明何时可把伽罗瓦群忠实嵌入对称群。

\section{阿廷固定域定理与基本定理的准备}
证明不同域嵌入的线性无关性，作为阿廷固定域定理（Artin's fixed-field theorem）的关键引理。若 $H$ 是域 $L$ 的一个有限自同构群，证明 $L/L^H$ 为有限伽罗瓦扩张且 $[L:L^H]=|H|$。与嵌入计数结合，解释有限伽罗瓦扩张的自同构数为何恰好等于扩张次数。

\section{有限伽罗瓦理论基本定理}
设 $L/F$ 为有限伽罗瓦扩张，$G=\operatorname{Gal}(L/F)$。建立中间域 $F\subseteq E\subseteq L$ 与子群 $H\leq G$ 之间的互逆对应
\begin{equation}\label{eq:galois-correspondence}
  E\longmapsto\operatorname{Gal}(L/E),
  \qquad H\longmapsto L^H.
\end{equation}
证明对应反转包含关系；当 $E=L^H$ 时，$[L:E]=|H|$、$[E:F]=[G:H]$。进一步证明 $E/F$ 为伽罗瓦扩张当且仅当 $H\trianglelefteq G$，此时限制映射诱导 $\operatorname{Gal}(E/F)\cong G/H$。区分总是伽罗瓦的 $L/E$ 与未必伽罗瓦的 $E/F$。

\section{伽罗瓦群的实际计算}
按“分裂域、扩张次数、根的可能像、关系式限制、子群判定”的次序组织计算。完整处理双二次扩张、$x^3-2$ 的分裂域与有限域扩张。对特征为 $0$ 的域上的不可约三次多项式引入判别式（discriminant），证明判别式是否为基域中的平方分别对应 $A_3$ 与 $S_3$。

\section{例题、练习与自检}
列出 $\mathbb Q(\sqrt2,\sqrt3)/\mathbb Q$ 的所有中间域，计算 $x^3-2$ 分裂域中各子群的固定域。练习要求验证某个根的置换确实保持全部代数关系，防止把任意置换都误判为自同构；另用非正规扩张说明不能直接套用式\eqref{eq:galois-correspondence}的全部结论。

